% Revize 2026-10-28: obsah, matematika a návaznost; podrobnosti a časy
% viz planovani/REVIZE_PRIPRAV_1_ROCNIK.md.
% Aktuální porovnání se skutečným obsahem místních PDF Realisticky.cz:
% externi_zdroje/realisticky/01_Zakladni_poznatky/06_Mocniny/010610_Mocniny_shrnuti/010610_Lekce.pdf
% externi_zdroje/realisticky/01_Zakladni_poznatky/04_Formalni_logika_vyroky/010410_Formalni_logika_shrnuti/010410_Lekce.pdf
% externi_zdroje/realisticky/04_Goniometrie/02_Goniometricke_funkce/040201_Goniometricke_funkce_ostreho_uhlu/040201_Lekce.pdf
% Hodina 32, 1. ročník SPŠ. Rozsah: Rozpis_vyuky_MAT_1_rocnik.tex, řádek 32.

% Plán: 1_rocnik/tematicke_plany/SPS.pdf; společná projektová šablona.

% Učebnice: ucebnice/01_Zakladni_poznatky/Mocniny_a_odmocniny/Mocniny a Odmocniny.tex

% Původ úloh: Vlastní zadání a vlastní formulace výkladu; metodická inspirace zdroji.

% Metodický zdroj: externi_zdroje/realisticky/01_Zakladni_poznatky/06_Mocniny/010610_Mocniny_shrnuti/010610_Lekce.pdf

% URL: http://www.realisticky.cz/ucebnice/01%20Matematika%20S%C5%A0/01%20Z%C3%A1kladn%C3%AD%20poznatky/06%20Mocniny/10%20Mocniny%20-%20shrnut%C3%AD.pdf

\documentclass[11pt,a4paper]{article}

% Spolecny styl pro ucitelske pripravy MAT (Work/MAT)
% Verze 1.1, 2026-10-23
\NeedsTeXFormat{LaTeX2e}
\ProvidesPackage{priprava}[2026/10/23 v1.1 Sablona ucitelskych priprav MAT]

% Zakladni typografie a cestina
\RequirePackage[T1]{fontenc}
\RequirePackage[utf8]{inputenc}
\RequirePackage[czech]{babel}
\RequirePackage{lmodern}
\RequirePackage{microtype}
\RequirePackage{amsmath,amssymb}
\RequirePackage{xcolor}
\RequirePackage{enumitem}
\RequirePackage[most]{tcolorbox}
\RequirePackage{tabularx,booktabs,array}
\RequirePackage{fancyhdr}
\RequirePackage{needspace}
\RequirePackage{geometry}
\RequirePackage[hidelinks]{hyperref}

\geometry{a4paper,top=17mm,bottom=18mm,left=23mm,right=23mm}
\definecolor{notegray}{gray}{0.39}
\setlength{\parindent}{0pt}
\setlength{\parskip}{2.5pt plus 1pt minus 0.5pt}
\setlist[itemize]{topsep=2pt,itemsep=1pt,parsep=0pt,leftmargin=1.45em}
\setlist[enumerate]{topsep=2pt,itemsep=2pt,parsep=0pt,leftmargin=1.65em}
\setcounter{secnumdepth}{0}
\renewcommand{\arraystretch}{1.12}

% Zahlavi: vlevo rocnik | globalni tema, vpravo cislo: nazev hodiny
% Zapati zustava zcela prazdne.
\pagestyle{fancy}
\fancyhf{}
\renewcommand{\headrulewidth}{0.35pt}
\renewcommand{\footrulewidth}{0pt}
\setlength{\headheight}{14pt}
\newcommand{\nastavPripravu}[4]{%
  \fancyhead[L]{\small #1\enspace\textbar\enspace #2}%
  \fancyhead[R]{\small #3: #4}%
}

% Minimalisticke cernobile boxy. Kazda definice zvlast.
\tcbset{commonlesson/.style={enhanced,breakable=false,colback=white,colframe=black,
  colbacktitle=white,coltitle=black,boxrule=0.45pt,
  arc=0.8mm,outer arc=0.8mm,left=3mm,right=3mm,top=1.6mm,bottom=1.6mm,
  toptitle=1mm,bottomtitle=0.7mm,fonttitle=\bfseries,
  attach boxed title to top left={xshift=2.5mm,yshift=-1.15mm},
  boxed title style={colback=white,colframe=white,boxrule=0pt,left=1mm,right=1mm}}}
\newtcolorbox{zadaniBox}[1][]{commonlesson,title={Zadání},#1}
\newtcolorbox{prikladBox}[1][]{commonlesson,title={Příklad},#1}
\newtcolorbox{definiceBox}[1][]{commonlesson,title={Definice},#1}

% Spolecny box pro zadani a jeho reseni, bez mezery mezi castmi.
% Pouziti:
%   \begin{prikladReseniBox}
%     Zadani ulohy.
%     \tcblower
%     \begin{align*} ... \end{align*}
%   \end{prikladReseniBox}
% Spodni cast ma jemny sedy podklad a delici linku; lze v ni sazet
% vice radku matematiky. Obe casti maji jen jeden spolecny ramecek.
\tcbset{solutionlesson/.style={
  commonlesson,
  enhanced jigsaw,
  breakable,
  colbacklower=black!3,
  segmentation style={solid,draw=black!25,line width=0.35pt},
  fontlower=\small\color{notegray},
  before lower={\noindent\textbf{Řešení:}\enspace},
  leftlower=3mm,rightlower=3mm
}}
\newtcolorbox{prikladReseniBox}[1][]{solutionlesson,title={Příklad},#1}
\newtcolorbox{zadaniReseniBox}[1][]{solutionlesson,title={Zadání},#1}

% Pomucky pro autora pripravy: sedive poznamky; \reseni pro starsi pripravy
\newcommand{\poznamka}[1]{\par\begingroup\small\color{notegray}\itshape
  \textnormal{\textbf{Poznámka pro učitele.}}\enspace #1\par\endgroup}
\newcommand{\reseni}[1]{\par\begingroup\small\color{notegray}
  \textbf{Očekávané řešení:}\enspace #1\par\endgroup}
\newcommand{\blok}[1]{\needspace{5\baselineskip}\subsection*{#1}}
\newcommand{\vyrok}[1]{\uv{#1}}


\nastavPripravu{1. ročník}{Opakování}{32}{Čtvrtletní písemná práce}

\begin{document}

\blok{Výklad a souvislosti}

Podle rozpisu je hodina určena souhrnné písemné práci z výrokové logiky, množin, goniometrie a mocnin. Tato příprava vede organizaci a následnou nápravu; neobsahuje žákovské zadání testu ani jeho varianty. Konkrétní test se připravuje zvlášť podle skillu \texttt{vytvor-test}.

\begin{definiceBox}[title={Připomenutí definice},breakable=false]
Pro $a\ne0$ je \textbf{mocnina se záporným exponentem} $a^{-n}=1/a^n$ při $n\in\mathbb N$. Definici při opravě používáme k vysvětlení chyby, nikoli jen k opsání pravidla.
\end{definiceBox}


Před hodinou připrav schválený test v uvedeném rozsahu a určuj organizační podmínky podle skutečného zadání učitele. Tato příprava sama nepovoluje kalkulačku ani jiné pomůcky. U souhrnné práce počítej s převážnou částí hodiny na samostatné psaní; nezkracuj ji kvůli povinnému rozboru všech oblastí. Následující příklady slouží jen rozboru po odevzdání nebo individuální zpětné vazbě po opravě. Pokud nezbude čas, odlož je na následnou konzultaci; neměň kvůli tomu rozpis dalších hodin.

Při opravě rozlišuj chybně přečtené zadání, nevhodný model, chybu ve znaménku a početní chybu. Zkontroluj také podmínky výrazů a jednotky geometrických výsledků. Veřejně nerozebírej výkony jednotlivých žáků.

\blok{Řešené příklady}

\begin{prikladReseniBox}[title={Příklad 1},breakable=false]
Pro následnou nápravu vysvětli rozdíl mezi $(-2)^{-2}$ a $-2^{-2}$.
\tcblower
$(-2)^{-2}=1/4$, ale $-2^{-2}=-1/4$. Nejprve určíme celý základ, pak použijeme převrácenou hodnotu.
\end{prikladReseniBox}


\begin{prikladReseniBox}[title={Příklad 2},breakable=false]
Při rozboru převeď $\{x\in\mathbb Z\mid -1\le x<3\}$ na výčet.
\tcblower
$\{-1,0,1,2\}$. Pravá mez je ostrá a nepatří do výčtu; necelá čísla vylučuje obor $\mathbb Z$.
\end{prikladReseniBox}


\blok{Samostatné procvičení}

\begin{zadaniReseniBox}[title={Procvičení},breakable=false]
\begin{enumerate}[label=\alph*)]
\item Po opravě neguj výrok „Každý prvek množiny $A$ je kladný“.
\item Při nápravě určete přeponu pravoúhlého trojúhelníku s odvěsnami $3\,\mathrm{cm}$ a $4\,\mathrm{cm}$ a sinus úhlu proti kratší odvěsně.
\item Po odevzdání testu vysvětli, proč $x^3:x^3=1$ vyžaduje $x\ne0$.
\end{enumerate}
\tcblower
\begin{enumerate}[label=\alph*)]
\item „Existuje alespoň jeden prvek množiny $A$, který není kladný“, tedy je nekladný. Nestačí říci, že všechny prvky jsou záporné.
\item Přepona je $5\,\mathrm{cm}$, protože $3^2+4^2=25$. Sinus uvedeného úhlu je $3/5$.
\item Pro $x=0$ bychom dělili nulou. Úprava na konstantu původní zákaz neruší.
\end{enumerate}
\end{zadaniReseniBox}


\poznamka{Úlohy v této přípravě nejsou návrhem dalšího testu. Vyber jen nápravu skutečně zjištěných chyb; zbytek je zásoba pro konzultaci. Na začátku ani během psaní nenapovídej matematické postupy. Bodování a klasifikační hranice zde nejsou zavedeny.}

\end{document}
